\paragraph{Blue~Gene/Q}
Our Blue~Gene/Q numbers were gathered on Vesta, a small Blue~Gene/Q
system located at Argonne National Laboratory.
Each compute node of the
Blue~Gene/Q is a 64-bit PowerPC A2 processor with 16 1.6 Ghz compute
cores and 16 GB of DRAM. The compute notes are grouped into drawers
(32 compute nodes per drawer), which are grouped into midplanes (16
compute drawers, 512 compute nodes).  Within the midplane, nodes are
electrically connected in a 5-D torus (4x4x4x4x2). Beyond the
midplane, an optical interconnect is used.  

All benchmarks were compiled using Native X10 version 2.4.0 with \texttt{-NO\_CHECKS}, \texttt{-O} options. The
generated C++ code was then compiled for Blue~Gene/Q using xlC v11.1
with
the \texttt{-O3} \texttt{-qinline} \texttt{-qhot} \texttt{-qtune=qp} \texttt{-qsimd=auto} \texttt{-qarch=qp}
command line arguments.

To maximize the number of X10 Places, all experiments with 16 Places or
more were run using the \texttt{c16} mode, which creates one MPI
process (i.e., one X10 Place) per Blue~Gene/Q compute core.  Thus each
X10 Place has 1 A2 core and 1 GB of DRAM available to it.  
